% Paper A: theory HELD until probe P1 decides the chemotherapy design.
% NOT \input by paper.tex. Drafted 2026-11-15; NOT checked by JS.
%
% Everything here serves the chemotherapy experiment (several constraints, one
% radius over the 96 test cohorts, the max-t fallback and Table 2's common
% scale), whose design P1 may change. To use: move the pieces wanted into
% paper_a_guarantee.tex (statements) and paper_a_proofs.tex (proofs), and the
% Section 4 text at the bottom into paper.tex. Labels and refs assume the
% main-text labels (thm:selection, prop:simulated, cor:multi, ...).

% ---------------------------------------------------------------------------
% Was Remark rem:simulated (iii).
(iii) $\mathcal{X}$ in the definition of $w$ may be any set that is fixed
independently of $\varepsilon$ and contains every decision the procedure can
return. When part of the decision is a fixed context, such as a cohort's
characteristics, and decisions are returned for contexts
$c_1,\dots,c_T$, $\mathcal{X}$ can be the union of their slices. One radius
then holds for all $T$ decisions at once: the probability that any of them
violates its clean-label constraint is at most $\epsilon$.

% ---------------------------------------------------------------------------
\paragraph{A joint radius for several constraints.}
Consider $m$ learned constraints $f_j(x;\theta_j(y_j)) \le b_j$, each
satisfying Assumption~\ref{ass:affine} with its own $s_j$ and $c_j$, with noise
$\varepsilon_j = \Sigma_j^{1/2}\xi_j$ satisfying Assumption~\ref{ass:noise},
and let $W_j = w_j(\xi_j)$ be the variable $W$ of
Theorem~\ref{thm:selection} for constraint $j$.

\begin{corollary}[Joint radius]\label{cor:joint}
Suppose the joint distribution of $(\xi_1,\dots,\xi_m)$ is known. Fix weights
$\omega_j > 0$ and set $R_j = \omega_j\hat R$, where $\hat R$ is computed as in
Proposition~\ref{prop:simulated} from upper bounds on $N$ simulated copies of
$W_\omega = \max_j W_j/\omega_j$. Then any $x^\star$ satisfying
$f_j(x^\star;\hat\theta_j) + R_j\|s_j(x^\star)\|_{\Sigma_j} \le b_j$ for all
$j$, possibly chosen as a function of the data, violates at least one
clean-label constraint with probability at most $\epsilon$.
\end{corollary}

\begin{proof}
If $x^\star$ violates constraint $j$, the argument of
Theorem~\ref{thm:selection} gives $W_j > R_j$, so $W_\omega > \hat R$. The
proof of Proposition~\ref{prop:simulated} applies verbatim to $W_\omega$, of
which the simulated values are independent copies.
\end{proof}

The joint radius is never worse than an even Bonferroni split with simulated
per-constraint radii, in the following sense. Let $q_j$ be the
$(1-\epsilon/m)$-quantile of $W_j$ and take $\omega_j = q_j$. Then
$\mathbb{P}(W_\omega > 1) \le \sum_j \mathbb{P}(W_j > q_j) \le \epsilon$, so
the $(1-\epsilon)$-quantile of $W_\omega$ is at most one and each
$R_j = \omega_j\hat R$ is, up to simulation error, at most $q_j$. It gains
most when the $W_j$ are strongly dependent, as when the constraints are
fitted on the same inputs. It needs the joint law of the noise across
constraints; Corollary~\ref{cor:multi} does not. On chemotherapy the five
toxicities are measured on the same patients, so that law is not known.

% ---------------------------------------------------------------------------
\paragraph{The certified level of any decision.}
For a decision $x$ and weights $\omega_j > 0$ (with $\omega_j = 1$ for a
single constraint), the certified radius is
\[
  \hat r(x) = \min_{j = 1,\dots,m}
    \frac{b_j - f_j(x;\hat\theta_j)}{\omega_j\,\|s_j(x)\|_{\Sigma_j}} ,
\]
where a term with $\|s_j(x)\|_{\Sigma_j} = 0$ is read as $+\infty$ if
$f_j(x;\hat\theta_j) \le b_j$ and as $-\infty$ otherwise. (When
$\|s_j(x)\|_{\Sigma_j} = 0$ the clean-label and fitted predictions at $x$
coincide almost surely, so the verdict on constraint $j$ is known.) A
decision $x$ satisfies the robust counterpart of every constraint at radii
$R_j = \omega_j t$ exactly when $\hat r(x) \ge t$. Let
$G(r) = \mathbb{P}(W_\omega > r)$, with $G(-\infty) = 1$ and
$G(+\infty) = 0$, and write $A$ for the event that $x^\star$ violates at least
one clean-label constraint.

\begin{proposition}[Certified level]\label{prop:certified}
Let $x^\star \in \mathcal{X}$ be chosen in any way from the data, including by
maximizing $\hat r$. Then
\begin{enumerate}[label=(\roman*)]
  \item $\mathbb{P}\bigl(A \text{ and } \hat r(x^\star) \ge r\bigr) \le G(r)$
    for every $r \in \mathbb{R}$;
  \item the level $\hat\epsilon = G(\hat r(x^\star))$ satisfies
    $\mathbb{P}(A \text{ and } \hat\epsilon \le \alpha) \le \alpha$ for every
    $\alpha \in [0,1]$.
\end{enumerate}
\end{proposition}

\begin{proof}
On $A$, let $j$ be a violated constraint. If $\|s_j(x^\star)\|_{\Sigma_j} = 0$
the $j$-th term of $\hat r(x^\star)$ is $-\infty$, so
$\hat r(x^\star) = -\infty$. Otherwise, as in Theorem~\ref{thm:selection},
\[
  W_j \ge \frac{-s_j(x^\star)^\top\varepsilon_j}{\|s_j(x^\star)\|_{\Sigma_j}}
  > \frac{b_j - f_j(x^\star;\hat\theta_j)}{\|s_j(x^\star)\|_{\Sigma_j}}
  \ge \omega_j \hat r(x^\star) .
\]
In both cases $W_\omega > \hat r(x^\star)$ on
$A$, which gives (i). For (ii), $G$ is nonincreasing and right-continuous, so
$\{r : G(r) \le \alpha\}$ is an interval $[r_\alpha, +\infty]$ with
$G(r_\alpha) \le \alpha$ when $r_\alpha$ is finite. On the event in (ii),
$G(\hat r(x^\star)) \le \alpha$ gives $\hat r(x^\star) \ge r_\alpha$, so (i)
at $r = r_\alpha$ bounds its probability by $\alpha$. The cases
$r_\alpha = \pm\infty$ are immediate. If $r_\alpha = +\infty$ the event
requires $\hat r(x^\star) = +\infty$, under which no constraint can be
violated. If $r_\alpha = -\infty$, then $\hat\epsilon = 1$ on
$\{\hat r(x^\star) = -\infty\}$, and on the rest of $A$,
$W_\omega > \hat r(x^\star) > -\infty$, an event of probability at most
$\lim_{r \to -\infty} G(r) \le \alpha$.
\end{proof}

$G$ is not known in closed form, but the simulated copies used for the radius
give a level with the same guarantee.

\begin{corollary}[Simulated level]\label{cor:sim-level}
Let $\overline W_1,\dots,\overline W_N$ be upper bounds on independent
simulated copies of $W_\omega$, as in Proposition~\ref{prop:simulated}. Then
\[
  \hat\epsilon_N = \frac{1 + \#\{i : \overline W_i \ge \hat r(x^\star)\}}{N+1}
\]
satisfies $\mathbb{P}(A \text{ and } \hat\epsilon_N \le \alpha) \le \alpha$
for every $\alpha \in [0,1]$, over the noise and the simulation jointly.
\end{corollary}

\begin{proof}
Let $W_1,\dots,W_N$ be the copies and
$p = (1 + \#\{i : W_i \ge W_\omega\})/(N+1)$. On $A$,
$W_\omega > \hat r(x^\star)$ by the proof of
Proposition~\ref{prop:certified}, so
$\{i : W_i \ge W_\omega\} \subseteq \{i : \overline W_i \ge \hat r(x^\star)\}$
and $p \le \hat\epsilon_N$. It remains to show
$\mathbb{P}(p \le \alpha) \le \alpha$. Let $K = \lfloor \alpha(N+1) \rfloor$.
The event $p \le \alpha$ says that at most $K - 1$ of the copies are at least
$W_\omega$. Order $W_\omega, W_1,\dots,W_N$ decreasingly, breaking ties
arbitrarily; the variable in position $t$ has at least $t - 1$ others at least
as large, so at most $K$ of the $N+1$ variables have at most $K-1$ others at
least as large. By exchangeability,
$\mathbb{P}(p \le \alpha) \le K/(N+1) \le \alpha$.
\end{proof}

Neither statement is conditional on the level reported. They do not say that
a decision reported at level $\alpha$ violates with probability at most
$\alpha$ given that it was so reported; they bound the probability that a
violating decision is reported at level $\alpha$ or better. Since $x^\star$
may be chosen to maximize $\hat r$, the decision with the highest certified
level carries a valid level even when the robust counterpart at the target
radius is infeasible.

% ---------------------------------------------------------------------------
% Section 4 text, held with the above (was the second half of
% "Certifying a decision, and the fallback").

The certified radius $\hat r(x)$ needs one evaluation of each fitted model and
of each $\|s_j(x)\|_{\Sigma_j}$ at $x$. It therefore applies to the decision
of any method, and Proposition~\ref{prop:certified} and
Corollary~\ref{cor:sim-level} put every method's decision on one scale.

When the robust counterpart at the target radius is infeasible, the decision
with the highest certified level solves $\max_{x \in \mathcal{X}} \hat r(x)$.
The set $\{x \in \mathcal{X} : \hat r(x) \ge t\}$ is the feasible set of the
robust counterpart at radii $R_j = \omega_j t$ and shrinks as $t$ grows, so
the maximum is found by bisection on $t$, each step a feasibility problem of
the same class as the robust counterpart. At the largest $t$ found, the
robust problem at that radius returns the cheapest decision with that level,
which is reported with its level $\hat\epsilon_N$
(Corollary~\ref{cor:sim-level}); the level is valid although the decision was
chosen to maximize $\hat r$.
